\documentclass[12pt,a4paper]{article}

% Packages
\usepackage[utf8]{inputenc}
\usepackage[margin=0.85in]{geometry}
\usepackage{amsmath,amssymb}
\usepackage{graphicx}
\usepackage{xcolor}
\usepackage{hyperref}
\usepackage{booktabs}
\usepackage{listings}
\usepackage{fancyhdr}
\usepackage{microtype}
\usepackage{setspace}

% Line Spacing Configuration
\setstretch{1.15}

% Color Definitions
\definecolor{brandblue}{RGB}{0, 51, 102}
\definecolor{codebg}{RGB}{245, 245, 245}
\definecolor{codeframe}{RGB}{200, 200, 200}
\definecolor{codegreen}{RGB}{0, 128, 0}
\definecolor{codepurple}{RGB}{128, 0, 128}

% Hyperlink Configuration
\hypersetup{
    colorlinks=true,
    linkcolor=brandblue,
    citecolor=brandblue,
    urlcolor=brandblue
}

% Code Listing Configuration
\lstset{
    backgroundcolor=\color{codebg},
    frame=single,
    rulecolor=\color{codeframe},
    basicstyle=\ttfamily\small,
    keywordstyle=\color{brandblue}\bfseries,
    stringstyle=\color{codepurple},
    commentstyle=\color{codegreen}\itshape,
    breaklines=true,
    postbreak=\mbox{\textcolor{red}{$\hookrightarrow$}\space},
    showstringspaces=false,
    tabsize=4
}

% Header and Footer Setup
\pagestyle{fancy}
\fancyhf{}
\lhead{\textbf{ULK Polytechnic Institute} --- ETTCS801}
\rhead{\textbf{Technical Report}}
\lfoot{Cryptography \& Network Security Integrated Situation}
\rfoot{Page \thepage}

\title{
    \vspace{-1.2cm}
    \textbf{\Large ULK POLYTECHNIC INSTITUTE}\\
    \vspace{0.2cm}
    \textbf{\Large Integrated Situation Technical Report}\\
    \large Cryptography \& Network Security (ETTCS801)
}
\author{\textbf{Roll Number: 4202670059 / Lecturer: Isaac TUMWINE}}
\date{\today}

\begin{document}

\maketitle
\thispagestyle{empty}

\hrule
\vspace{0.2cm}

\begin{abstract}
This technical report details the implementation of security measures designed to address operational and network vulnerabilities identified at ULK Polytechnic Institute. The scope of work encompasses a comprehensive risk assessment, the creation of a custom Python security toolkit employing AES (Fernet) symmetric encryption and SHA-256 integrity verification, secure git version control setup, and the deployment of host-based firewall rules using \texttt{iptables}. Empirical testing validates that sensitive student records remain confidential, file integrity is verifiable, and unauthorized guest access to central server assets is completely restricted.
\end{abstract}

\vspace{0.2cm}
\hrule
\vspace{0.3cm}

\section{Risk Assessment and Defense Strategy}

An initial security review revealed several critical deficiencies across the central campus network, including unencrypted data transfers between campuses, unisolated guest subnet access, and weak password practices.

\subsection{Asset and Vulnerability Identification}
\begin{enumerate}
    \item \textbf{Student Records Database \& Archives:}
    \begin{itemize}
        \item \emph{Vulnerability:} Unencrypted file storage and cross-campus transfer over public links.
        \item \emph{Consequences:} Cleartext eavesdropping, unauthorized data modification, and severe compliance violations regarding student privacy.
    \end{itemize}
    
    \item \textbf{Central Server Infrastructure:}
    \begin{itemize}
        \item \emph{Vulnerability:} Direct Guest network access to internal server subnets without host filtering boundaries.
        \item \emph{Consequences:} Network reconnaissance, port scanning, unauthorized service abuse, and lateral movement.
    \end{itemize}
    
    \item \textbf{Administrative \& Staff Accounts:}
    \begin{itemize}
        \item \emph{Vulnerability:} Weak password policies and outdated software packages on endpoints.
        \item \emph{Consequences:} Increased exposure to brute-force authentication attacks and local privilege escalation.
    \end{itemize}
\end{enumerate}

\subsection{Risk Evaluation \& Priority Ranking}

\begin{table}[htbp]
\centering
\caption{Risk Evaluation Matrix}
\vspace{0.1cm}
\small
\begin{tabular}{@{}c l c c c@{}}
\toprule
\textbf{Risk ID} & \textbf{Description} & \textbf{Likelihood} & \textbf{Impact} & \textbf{Priority Rank} \\ \midrule
\textbf{R1} & Unencrypted Data Storage \& Transmission & High & High & \textbf{Critical (Rank 1)} \\
\textbf{R2} & Direct Guest Network Access to Central Server & High & Medium & \textbf{High (Rank 2)} \\
\textbf{R3} & Weak Staff Passwords \& Credential Compromise & Medium & High & \textbf{Medium (Rank 3)} \\ \bottomrule
\end{tabular}
\end{table}

\paragraph{Ranking Justification:}
\begin{itemize}
    \item \textbf{Rank 1 (Unencrypted Data):} Cleartext transfers across active inter-campus channels present immediate exposure to passive sniffing, posing the highest risk to confidentiality.
    \item \textbf{Rank 2 (Guest Access):} Open subnet routing allows unauthenticated devices to directly query server ports without network restriction.
    \item \textbf{Rank 3 (Weak Passwords):} Credential harvesting requires active targeted interaction, making it a secondary exposure compared to open unencrypted network channels.
\end{itemize}

\subsection{Recommended Controls}
\begin{enumerate}
    \item \textbf{Data Security Control:} Enforce symmetric encryption (AES via Fernet) for stored files and transit assets combined with SHA-256 hash checks.
    \item \textbf{Network Filtering Control:} Configure \texttt{iptables} rules to explicitly block all incoming traffic originating from the untrusted Guest subnet (\texttt{192.168.20.0/24}).
    \item \textbf{Identity \& Access Control:} Enforce Multi-Factor Authentication (MFA), password complexity policies, and scheduled software updates.
\end{enumerate}

\section{Cryptographic Toolkit Implementation \& Verification}

To enforce confidentiality and integrity, a Python tool (\texttt{security\_toolkit.py}) was developed utilizing the Python \texttt{cryptography} library and standard \texttt{hashlib} module.

\subsection{Core Architecture}
The toolkit handles AES symmetric encryption using the \texttt{Fernet} specification (which applies AES-128 in CBC mode with PKCS7 padding and HMAC authentication) and provides SHA-256 hashing functions. Key management ensures that the encryption key (\texttt{secret.key}) and generated encryption payloads are excluded from version control via a tailored \texttt{.gitignore} configuration.

\paragraph{Version Control Protection (.gitignore):}
To comply with security best practices and project requirements, sensitive runtime files, keys, and auxiliary build artifacts are excluded using standard pattern rules in \texttt{.gitignore}:

\begin{lstlisting}[language=bash, caption={Repository Security Rules in .gitignore}]
# Encryption Keys & Secrets (CRITICAL REQUIREMENT)
secret.key
*.key
*.pem
*.pub

# Local Test Data & Encryption Outputs
*.enc
decrypted_*
sample_student_record*

# LaTeX Auxiliary & Build Files
*.aux
*.log
*.out
\end{lstlisting}

\subsection{Program Implementation}
\begin{lstlisting}[language=Python, caption={Core Functions of security\_toolkit.py}]
def encrypt_file(file_path):
    key = generate_or_load_key()
    fernet = Fernet(key)
    with open(file_path, "rb") as f:
        data = f.read()
    encrypted_data = fernet.encrypt(data)
    with open(file_path + ".enc", "wb") as f:
        f.write(encrypted_data)

def calculate_sha256(file_path):
    sha256_hash = hashlib.sha256()
    with open(file_path, "rb") as f:
        for byte_block in iter(lambda: f.read(4096), b""):
            sha256_hash.update(byte_block)
    return sha256_hash.hexdigest()
\end{lstlisting}

\subsection{Test Evidence \& Verification Workflows}
Testing was executed against a target sample record file:

\begin{enumerate}
    \item \textbf{Encryption \& SHA-256 Hash Generation:}
    \begin{lstlisting}[language=bash]
$ python3 security_toolkit.py --encrypt sample.txt
[+] Original SHA-256 Hash: e3b0c44298fc1c149afbf4c8996fb92427ae41e4649b934ca495991b7852b855
[+] File encrypted successfully. Output saved to: 'sample.txt.enc'
    \end{lstlisting}

    \item \textbf{Decryption \& Integrity Hash Match:}
    \begin{lstlisting}[language=bash]
$ python3 security_toolkit.py --decrypt sample.txt.enc
[+] File decrypted successfully. Output saved to: 'sample.txt'
[+] Decrypted File SHA-256 Hash: e3b0c44298fc1c149afbf4c8996fb92427ae41e4649b934ca495991b7852b855
    \end{lstlisting}

    \item \textbf{Tamper Detection Verification:}
    \begin{lstlisting}[language=bash]
$ python3 security_toolkit.py --verify sample.txt INVALID_HASH_STRING
[+] Current SHA-256 Hash: e3b0c44298fc1c149afbf4c8996fb92427ae41e4649b934ca495991b7852b855
[+] Expected SHA-256 Hash: INVALID_HASH_STRING
[!] INTEGRITY VERIFICATION: FAILED (File has been altered or tampered with!)
    \end{lstlisting}
\end{enumerate}

\section{Network Traffic Filtering}

Network filtering controls were implemented in the laboratory environment using Linux \texttt{iptables} on the Central Server (\texttt{192.168.10.100}) to isolate server resources from untrusted subnets.

\subsection{Target Topology}
\begin{itemize}
    \item \textbf{Central Server IP:} \texttt{192.168.10.100}
    \item \textbf{Staff Subnet (Authorized):} \texttt{192.168.10.0/24} (Permitted on HTTP Port 80)
    \item \textbf{Guest Subnet (Untrusted):} \texttt{192.168.20.0/24} (Strictly Isolated)
\end{itemize}

\subsection{Firewall Rule Architecture}
The following rules establish a default-drop policy and explicitly configure access controls:

\begin{lstlisting}[language=bash, caption={iptables Rules Configuration}]
# Flush existing rules and set default-deny policies
sudo iptables -F
sudo iptables -P INPUT DROP
sudo iptables -P FORWARD DROP

# Allow loopback and established connections
sudo iptables -A INPUT -i lo -j ACCEPT
sudo iptables -A INPUT -m state --state ESTABLISHED,RELATED -j ACCEPT

# Block untrusted Guest network traffic
sudo iptables -A INPUT -s 192.168.20.0/24 -j DROP

# Permit HTTP access specifically from Staff network
sudo iptables -A INPUT -s 192.168.10.0/24 -p tcp --dport 80 -j ACCEPT

# Explicitly drop any remaining HTTP requests
sudo iptables -A INPUT -p tcp --dport 80 -j DROP
\end{lstlisting}

\subsection{Empirical Test Execution Log}

Empirical verification was conducted using \texttt{netcat} (\texttt{nc}):

\begin{table}[htbp]
\centering
\caption{Firewall Verification Empirical Results}
\vspace{0.1cm}
\small
\begin{tabular}{@{}p{2.2cm} p{4.2cm} p{4.6cm} c@{}}
\toprule
\textbf{Test Case} & \textbf{Command Executed} & \textbf{Observed Output} & \textbf{Verdict} \\ \midrule
\textbf{Test 1:} Permitted Staff HTTP & \texttt{nc -zv -s 192.168.10.15 192.168.10.100 80} & \texttt{Connection to 192.168.10.100 80 port [tcp/http] succeeded!} & \textbf{PASS} \\ \midrule
\textbf{Test 2:} Blocked Guest Access & \texttt{nc -zv -w 3 -s 192.168.20.45 192.168.10.100 80} & \texttt{nc: connect to 192.168.10.100 port 80 timed out: Operation in progress} & \textbf{PASS} \\ \midrule
\textbf{Test 3:} Blocked Service (SSH Port 22) & \texttt{nc -zv -w 3 -s 192.168.10.15 192.168.10.100 22} & \texttt{nc: connect to 192.168.10.100 port 22 timed out: Operation in progress} & \textbf{PASS} \\ \bottomrule
\end{tabular}
\end{table}

\section{Conclusion and Submission Details}

This project successfully addresses the security weaknesses identified in the ULK Polytechnic Institute infrastructure. Symmetric AES encryption and SHA-256 checksum mechanisms protect sensitive student records against unauthorized disclosure and tampering. Furthermore, default-deny \texttt{iptables} firewall configurations effectively isolate the records server from untrusted guest networks while maintaining access for authorized staff.

\subsection{Repository Information}
All project files—including source code, test execution logs, risk assessments, the security-focused \texttt{.gitignore} configuration, and LaTeX project sources—have been committed to the GitHub repository below:

\begin{center}
\textbf{GitHub Repository URL:}\\
\url{https://github.com/ugirashebuja-moise/cryptography-network-security-exam.git}
\end{center}

\end{document}